\section{Observability, cache provenance and bottleneck limits}\label{obs-observability}
The following observations concern the synthetic campaign, not Sui replay. Primary cases and the five separately executed 1,024-byte concurrent-write diagnostic cases are distinct populations. Each case has one repetition. The retained machine-readable \path{data/observability-evidence.json} includes exact case identifiers, raw stage counters, read-boundary deltas, host device and cgroup intervals, replay queries and units. Source is campaign 40848651463, at revision \texttt{ba3a6a0194cca6f171e5}\allowbreak\texttt{916649e7f4eeba65aeb0}. Stage definitions at that revision are in the following source files: \path{benchmarks/kv-research/src/backends/tidehunter.rs}; \path{tidehunter/src/distributed/tcp/diagnostics.rs}.

\subsection{Coverage and offline verification}\label{obs-coverage}
The monitoring capture reports 258 expected cases, 258 seen and 258 terminal observations, with no capture gaps under its own completeness criterion. It retained 378 query page receipts across 27 metric families, 339 unique page files, and a TSDB block with 4,508,934 samples and 195,976 series. Host and explicitly selected cgroup targets were scraped every 5 seconds, benchmark gauges every second. Initial verification reported 13 targets and 27 executable panels; the final Cassandra-anchored panels include empty queries. Missing data is not zero. The collector was cohosted and its overhead was not calibrated. Short cases and rolling rate windows need not align: host-wide CPU, memory, network, pressure and device signals include background and observation work. Selected cgroups are not a complete process-discovery or cross-backend cache-coverage guarantee.

Using the retained archive and already cached pinned images, an offline replay restored original timestamps; authenticated Grafana/datasource verification completed successfully and the replayed Prometheus host and resource queries returned data. Only the two owned replay containers were stopped afterwards. This was API/query verification, not visual inspection of Grafana. At Unix time 1793270529.054, near the capture end, all six host series reported swap total and swap free as zero, and the \texttt{nvme1n1} in-flight I/O gauge was zero at that instant. Host available memory was 119.77--121.00 billion bytes, and file-cache gauges were 0.359--1.558 billion bytes. These are host-level snapshots, not cgroup memory usage, peak disk queue depth during the cases, or evidence of absent stalls. A 20-second non-loopback receive-rate query returned points; it mixes host/background traffic and is not a per-case network accounting. ENA and Linux network-device counters can omit OS-bypass EFA traffic: absent or low counters do not establish an idle EFA NIC.

\subsection{What measured write stages can localize}\label{obs-stages}
Each distributed diagnostic captures the first 256 append requests on each of five append-owning connections after measurement arming: 1,280 selected requests, not a random sample and not the 4,096 logical writes in its coordinator interval. Client exchange, server handler, owner execution, persistence and transport clocks overlap or nest. Stage sums across requests and machines are not elapsed case time; queue residence through local completion is not peer receipt. The connection receive/idle field is not attributable under multiplexing and is retained as zero, not measured zero idle. The coordinator baseline is taken at arm time and the end before profile finalization. The following sums are descriptive diagnostic evidence in milliseconds; they are neither latency percentiles nor an additive critical-path decomposition.

\begin{center}\small
\begin{tabular}{lrrrrrr}
\hline
Diagnostic & Owner append & Append calls & Persist & Persist calls & WAL file & File calls\\
\hline
TCP & 156.319 & 1280 & 217.045 & 1280 & 165.336 & 1099\\
EFA & 120.249 & 1280 & 203.480 & 1280 & 165.815 & 1165\\
EFA direct & 110.205 & 1280 & 194.981 & 1280 & 163.056 & 1182\\
EFA native label & 117.908 & 1280 & 199.573 & 1280 & 165.096 & 1160\\
\hline
\end{tabular}
\end{center}
For example, the EFA-native diagnostic recorded 468.276 ms client exchange summed across 1,295 client requests, and 74.787 ms control-queue residence across 1,295 queue samples. These denominators differ from the 1,280 owner append calls and 1,160 file-persist calls, so dividing every clock by 1,280 or comparing the sums as disjoint components is invalid. Its coordinator interval records 4,096 admitted and completed writes, 2,207 append requests, 1,854 durable-receipt publication-blocked batches and one dependency-wait batch; these counters are not per-call tail attribution. The standalone native diagnostic selects the first 256 scalar writes: append sum 103.895 ms, index application 0.890 ms, and separately sampled WAL file persistence 176.219 ms across 512 file calls and namespace persistence 3.702 ms across 256 calls. Persistence components are nested in the foreground barrier, not extra stages to add to it. Batch native writes and per-operation stage percentiles are unsupported.

Current primary 1 KiB write p99 is 8.765 ms TCP, 14.881 ms EFA, 14.437 ms direct and 15.074 ms EFA-native, versus 3.311 ms native. Each write population has 4,096 calls, fewer than 10,000; its p99 is therefore flagged as sparsely supported and is not a run-to-run uncertainty estimate. Their historical counterparts were 144.163, 142.793, 142.884, 142.371 and 3.318 ms respectively. The older Frankfurt fleet and current Virginia fleet differ: this is descriptive, not a controlled code-only comparison. Current summed diagnostic clocks and interval counters show work in several stages but do not reveal which stage caused an individual primary p99, much less isolate the old tail. No one-way network duration follows from subtracting unsynchronized node clocks or from adding/subtracting independently computed quantiles.

\subsection{Read-path provenance and physical-I/O boundaries}\label{obs-cache}
The primary Tidehunter configuration disables explicit coordinator and owner value caches; it does not make the Linux page cache cold. In the EFA-native primary 1 KiB concurrent-read case, measurement-boundary deltas show 8,192 coordinator cache bypasses, 8,162 remote read RPCs and 8,162 logical owner WAL reads and async submissions. They do not all have a one-to-one relation to the 8,192 calls: coordinator misses total 15,516. Async submission and logical WAL read are not proof of physical-device access. Owner snapshots are sequential and not globally atomic.

Configured-cache cohorts are separate from primary. In case 0078 (EFA-native, 256 B, concurrency 1), the measurement boundary shows 8,192 coordinator-cache hits, zero misses and zero remote RPCs; API p50/p99 are 0.870/1.810 microseconds. This is a local cache-hit workload, not evidence of remote 1.810 microsecond p99 and not directly comparable to native at concurrency 16. In all eight owner-cache cases (256 B and 8 KiB across four transports), 8,192 owner lookups comprise 5,435 owner-cache hits and 2,757 misses/logical WAL reads. This explicit owner value cache is different from the Linux resident-WAL/NOWAIT path and from the kernel page cache. The 8 KiB TCP coordinator-cache case has 8,188 hits and four misses/RPCs; the other seven coordinator-cache cases have 8,192 hits and zero misses/RPCs.

Registered-RAM EFA-native cases are mixed populations: the 256 B case has 2,372 RMA hits/completions and 5,820 message requests, and the 8 KiB case 2,361 and 5,831 respectively, each out of 8,192, with zero RMA errors. The corresponding owner snapshots also report 5,820 and 5,831 successful resident reads. Thus roughly 29\% of calls used RMA in these cases; aggregate latency is not pure RMA latency. The EFA-native label does not imply all primary reads are one-sided.

All nine 64 KiB cold-OS-cache cases record a requested OS cache drop on all six hosts but \texttt{application\_\allowbreak caches\_\allowbreak evicted=false}; this cannot evict all database caches. The whole-host instance-store device read deltas support physical activity during boundaries, not attribution to individual GETs. For example EFA-native storage-0 reports 139,856 sectors, or 71,606,272 bytes, while its selected storage cgroup reports 70,995,968 read bytes; native's coordinator reports 546,104 sectors, or 279,605,248 bytes. These scopes and sampling boundaries differ. Valkey's dataset remains RAM-resident despite host I/O; FoundationDB, TiKV and Cassandra may retain their own caches. No common nine-backend hit ratio or universal per-request physical-NVMe claim is established.

\subsection{Resource signals and cautious interpretation}\label{obs-resource}
Matched primary 1 KiB concurrent reads show EFA-native storage-0 using 6.302326 CPU seconds over 1.556699393 seconds, about 4.049 average cores, and zero observed throttled periods in that selected cgroup interval; its other four storage roles average 4.046--4.052 cores. The TCP storage roles average 0.126--0.135 cores in their own intervals. EFA-native's client reports one throttled period of nine, with 0.022251 seconds throttled; the selected storage cgroups report no throttled periods. These intervals include gating/control and background activity and are not CPU per operation, instantaneous saturation, or proof that polling caused a latency difference. All selected EFA-native storage cgroups report zero read bytes in that warm boundary interval; that cgroup accounting does not prove every read was cached or identify physical reads outside its scope. Cgroup memory current is an end-boundary gauge, while memory peak is a cgroup-lifetime high-water mark, not interval RSS. Workload ceilings were 8 vCPU and 16 GiB per server/client host, with role-specific allocations possible.

The read-path deltas, CPU contrast, persistence and queue sums provide hypotheses for targeted instrumentation; none independently establishes a causal bottleneck. Additional independent repetitions, correlated per-request spans with synchronized clocks, verified EFA provider/NIC counters, physical-I/O attribution and calibrated observer overhead would be needed to localize tails more confidently.
